%% 毕业设计小结

\begin{designsummary}

    本次毕业设计以载人飞船发射逃逸系统作为研究对象，面向载人登月的任务需求，完成了六自由度动力学建模、蒙特卡洛打靶仿真与风险评估。经过本课题的训练，我在航天动力学建模、数值仿真和数据分析等方面得到了较为全面的锻炼。

    课题推进过程中遇到了一些困难。一是数据查找环节，猎户座飞船作为美国重要的载人飞船，其气动参数、发动机推力数据和惯量参数等详细技术指标并不公开，一些关键参数需要通过估算获得，一定程度上影响了仿真精度，但也让我积累了在不完整信息条件下开展工程分析的实践经验。二是干扰项分布参数的确定，一些干扰项分布及参数并不明确，在反复查阅文献和与老师讨论后才确定了最终的参数方案。三是在引入风扰进行零高度逃逸打靶时，没有考虑姿态奇异角保护，导致部分弹道的姿态角发散，最终仿真结果失真，后续通过对姿态角做限幅处理解决了该问题，这一经历让我深刻体会到数值仿真中数值稳定性保护机制的重要性。

    本课题的研究工作锻炼了我对多门专业课程知识的综合运用。六自由度运动方程推导和坐标系变换是航天飞行动力学课程的基础框架；姿态控制模块中的比例微分控制律是自动控制原理课程的重点之一；正交试验设计、拉丁超立方抽样和敏感性分析的统计学支撑则来自于概率论与数理统计课程。三门课程的知识在本课题中形成了完整的方法论链条。

    通过本次毕业设计，我加深了对航天器逃逸系统的理解，掌握了六自由度仿真建模和蒙特卡洛方法的工程应用，文献调研、独立分析与问题排查等科研能力上，也有了明显进步。未来，我将继续关注航天技术的发展，努力学习理论知识，争取为相关领域的研究与应用贡献自己的力量。

\end{designsummary}
